\documentclass[11pt]{amsart}
\input{../commands.tex}

\newtheorem{hwexercise}{Exercise}

\title{math231br Homework 6 (Due Thursday, 3/12/26)}

\begin{document}
\maketitle


\begin{hwexercise} This is basically just checking definitions, but argue carefully that a vector bundle is an example of a fiber bundle, with structure group $\GL_n(\R)$ and fiber $\R^n$.
\end{hwexercise}

\begin{hwexercise} Let $T^2 \to S^1$ be the torus, viewed as an $S^1$-bundle over $S^1$, with structure group $C_2$ given by rotating the circles in the fibers 180 degrees. Let $T^2_\text{twist} \to S^1$ be a similar fiber bundle, but where we cut the original one, turn one end of it 180 degrees and reglue it.
\begin{enumerate}
    \item Are $T^2$ and $T^2_\text{twist}$ homeomorphic as $S^1$-fiber bundles with structure group $C_2$ over $S^1$?
    \item If not, can you enlarge the structure group so that they become homeomorphic?
\end{enumerate}
\end{hwexercise}

\begin{hwexercise} Try your very hardest to define higher nonabelian \v{C}ech cohomology $\check{H}^n(X,G)$ when $G$ is non abelian. Then declare defeat and give up.
\end{hwexercise}

\begin{hwexercise} Let $G$ be a topological group acting transitively on a Hausdorff space $Y$. Fix any point $y\in Y$, and denote by $G_y$ the isotropy subgroup
\begin{align*}
    G_y = \left\{ g\in G \mid gy =y \right\}.
\end{align*}
Prove that $G_y \le G$ is a closed subspace.
\end{hwexercise}

\begin{hwexercise} Let $G$ be a topological group and $H \le G$ a closed subgroup. Prove that $G/H$ is Hausdorff.
\end{hwexercise}

\begin{hwexercise} Prove that a principal $G$-bundle is trivial if and only if it admits a section.
\end{hwexercise}





\end{document}
